%======================================================================
\section{Reproducibility}
\label{app:reproducibility}
%======================================================================

\subsection{Code, results and provenance}

All results come from one Python package and its scripts. The package
(\texttt{ks\_latent}) contains the PDE solvers, the autoencoders and
propagators, the training loops including window training, the
Lyapunov analysis, and the filters; scripts run each experiment and write
a JSON file per run with its configuration, seeds and wall time. Every
table body and figure in this paper is produced by a script in
\texttt{scripts/paper\_figures/} that reads only these result files; each
figure is written with a sidecar file listing its inputs, the commit and a
time stamp. A decision log records every choice of truth, model and
convention and every correction made while writing. A number-provenance
check compares every decimal number in the text with the generated
tables and the result files and lists those it cannot match, each of
which has been checked by hand.
% TODO(G8): replace by the tagged release commit and archive DOI once the code is released; the
% results were produced from a working tree on top of commit 2cd9da0 with uncommitted changes.

\subsection{Numerics, devices and seeds}

The PDEs are integrated in double precision on the CPU with ETDRK4. The
networks are trained in single precision on the GPU of an Apple M5 Pro
(48~GB unified memory); training there is not bitwise reproducible, so
repeated trainings with one seed agree only statistically, and we report
separate training seeds where we have them. Lyapunov spectra propagate
tangent vectors in single precision and re-orthonormalize them in double
precision on the CPU (Appendix~\ref{app:model_spectra}); all filter linear
algebra is done in double precision on the CPU. Initial conditions are
indexed by seed (0, 1, 2) and spun up for 100--500 time units. PDE
trajectories from one seed are reproducible only for a fixed number of
FFT threads; runs that differ in thread count follow different
trajectories after a few Lyapunov times, and we use them as independent
realizations (Appendix~\ref{app:groundtruth}). The software environment
is Python 3.11.15, PyTorch 2.13.0, NumPy 1.26.4 and SciPy 1.17.1, fixed by
an exported environment file.

\subsection{Tests}

The properties the paper relies on are checked by unit tests, 132 of which
cover the claims below and pass in a few seconds:
\begin{itemize}
\item \emph{Solver:} the linear growth spectrum, the exponential
  time-differencing coefficients against high-precision evaluation, mean
  conservation, translation equivariance and fourth-order convergence.
\item \emph{Architecture:} exact translation equivariance of the 2D
  autoencoder along each axis (Proposition~\ref{prop:equivariance}); the
  gauge-anchored channel equals the patch average; the decoder Jacobian is
  exactly block-diagonal at zero reach (Proposition~\ref{prop:blockdiag}).
\item \emph{Window training:} a forced-halo rollout reproduces the
  interior of the same propagator run on the large periodic domain, in 1D
  and 2D, and a halo one site too narrow breaks the equality
  (Proposition~\ref{prop:halo}).
\item \emph{Lyapunov analysis:} the double-precision tangent path
  recovers the spectrum of a linear map to machine precision.
\item \emph{Observation Jacobian:} the coloured construction equals
  reverse-mode differentiation exactly, and the observation-to-site map
  is checked pixel by pixel; deployment grids keep the training spacing.
\item \emph{Filters:} with a linear observation operator and Gaussian
  prior, the particle flow filter, the dense and per-site ensemble Kalman
  filters and the serial filter converge to the Kalman posterior at large
  ensemble size; the deterministic per-site transform reproduces the
  Kalman update with the ensemble covariance exactly; sites beyond the
  taper support are left unchanged; the projected and serial forms of
  sampling error correction reduce to the plain filter for an identity
  correction.
\end{itemize}

\subsection{Compute}

All runs used one workstation. Typical costs: a 1D Lyapunov spectrum of
the PDE at $L=400$, 35~min; a 120-unit Kolmogorov truth at $3\times$, several
hours; a model spectrum at $3\times$ or $88^2$, about 7~min; training a 2D
propagator, 17~min to 2~h; a data assimilation run of 100 cycles,
$9$--$97$~s of wall time up to $88^2$ and $4$--$59$~min at $220^2$--$440^2$, with one or four
tilings (Table~\ref{tab:da_cost}).

\subsection{Part~II code}

The Part~II experiments are run by \texttt{scripts/exp\_h9\_tiled\_da.py} (the latent filter with
shifted tilings, forecasts and the robustness options), \texttt{scripts/exp\_i11\_pde\_letkf.py} (the
physics-based filter with wrong models, parameter estimation, bias correction and the hidden-field
truth), \texttt{scripts/exp\_i24\_phys\_emulator.py} (the learned model inside a physical-space
filter), \texttt{scripts/exp\_j16\_estimate\_pde\_model\_error.py} (offline one-cycle residuals),
\texttt{scripts/exp\_j17h\_fit.py} (best fits of the wrong-damping models) and
\texttt{scripts/exp\_j18b\_roundtrip.py} (the round-trip test). The noisy-data model is trained by
\texttt{scripts/train\_kolmogorov\_crop.py --data-noise 0.1}. The library modules are
\texttt{ks\_latent/da/letkf\_local.py} (per-site filter with RTPS), \texttt{letkf\_phys.py}
(physical-space filter, multi-field states), \texttt{param\_est.py} (parameter estimation),
\texttt{model\_error.py}, and in \texttt{ks\_latent/solver/} the wrong-model operators of
\texttt{ks2d.py}, the cached imperfect-model truths of \texttt{twin.py} and the hidden-field system of
\texttt{ks2d\_memory.py}, each with unit tests (the hidden-field solver against the exact solution of
its linearization; the joint parameter update against the ensemble Kalman regression). The queues
that ran every Part~II experiment are kept as \texttt{scripts/run\_*.sh}.
